\section{总结、已知局限与后续演进}

\subsection{项目交付总结}

回到第 1 章提出的自由行规划四大痛点，本项目围绕“\textbf{可核对、可微调、本地可控、易安装维护}”四项原则，完成了从需求收敛、架构决策、核心引擎研发到多维度实测验收的完整工程闭环：

\begin{itemize}
  \item \textbf{架构与工程扎实}：以单应用容器统一交付 React + TypeScript 前端与 FastAPI + SQLite (FTS5, Schema v6) 后端，配套进程独占锁、事务迁移与 \texttt{travel-data} 离线校验/备份/空卷恢复工具，同时支持 Docker Compose 一键部署与 \texttt{git clone} 本机运行。
  \item \textbf{规划与核验可信}：动态 LLM 动作循环配合程序规则校验、独立质量审核与检查点续跑，彻底告别了“一错全废”或“盲目信任大模型”的两难；分离式三段原文引用与 \texttt{Exa $\rightarrow$ DeepSeek $\rightarrow$ Tavily $\rightarrow$ AnySearch} 四家搜索回退，使营业时间、淡旺季票价与预约要求真正做到有据可查、未知不冒充通过。
  \item \textbf{体验克制且验证充分}：自由用餐时段、纯净活动介绍、定点局部编辑与只读历史版本让行程真正具备可操作性；349 项后端测试、60 项前端测试、28 组 Playwright 流程（38 个状态 WCAG A/AA 0 违规）、西安五日游 0 硬冲突实测、个人偏好 A/B/C 真实验证与 Linux 容器空卷恢复，共同构成当前版本的坚实交付证据。
\end{itemize}

\subsection{当前已知局限}

基于本轮真实在线验收与压测数据，系统仍存在四类客观边界，值得在使用与后续迭代中坦诚面对：

\begin{enumerate}
  \item \textbf{长上下文多轮规划耗时偏长}：在外部模型串行决策与取证审核模式下，真实五日游规划与局部编辑耗时分别达到约 26 分钟与 4.3 分钟，期间偶发单步 \texttt{provider\_timeout}（需依赖检查点重试与降级机制保底，因此长行程建议将 \texttt{PROVIDER\_TIMEOUT\_SECONDS} 设为 120 秒）。
  \item \textbf{官方事实覆盖仍受公开搜索生态制约}：虽然故宫、兵马俑、碑林等核心场馆已能核实部分官方淡季票价与预约要求，但公开搜索引擎返回结果仍以第三方攻略和聚合页居多，导致行程中仍有相当比例的营业时间与费用停留在 \texttt{estimated} 或 \texttt{unknown}，最终呈现为 \texttt{needs\_review}。
  \item \textbf{复杂远郊跨区排程偶发模型建模失误}：如个人偏好 A 组对照实验所示，当行程跨市区与远郊（钟楼至兵马俑 41.2 km）且伴随“短步行”偏好时，模型偶发仍可能写出不合理的活动间隔与交通方式；虽然程序规则校验已可靠拦截为 \texttt{blocked} 避免误导用户，但仍需进一步提升草案自修复成功率。
  \item \textbf{交互与功能边界留白}：当前偏好更新采用用户在设置页主动编辑保存的方式，尚未在行程编辑后自动弹出“建议写入长期偏好”的确认条；天气预报、多城串联行程以及扫描件 PDF 的本地 OCR 均按首版简化原则留待后续迭代。
\end{enumerate}

\subsection{后续演进方向}

面向下一阶段的实用化打磨，建议按投入产出比依次推进以下四项工作：

\begin{itemize}
  \item \textbf{取证并发与官方站点定向召回}：在遵守各搜索提供方 QPS 与冷却机制的前提下，对独立场馆的事实查询引入受控并发并强化对官网/政府域名（如 \texttt{.gov.cn}、景区官网）的优先检索策略，同时缩短整体规划时长并提升 \texttt{supported} 事实占比。
  \item \textbf{远郊与跨区交通规则预检}：在草案生成后、消耗修复轮次前增加显式的长距离步行纠偏提示（例如超过 3 公里自动回传公交/驾车耗时对照给模型），减少类似 A 组实验中的远郊交通冲突。
  \item \textbf{偏好提炼确认交互}：在局部编辑完成后可选异步提炼一条简短的长期偏好更新建议，以非侵入式提示条展示在页面，由用户一键确认或忽略，进一步打通 R03 的体验闭环。
  \item \textbf{可选天气与真实用户试用反馈}：按既有接口预留接入和风天气 7--10 天预报作为室内外排程参考，并在个人真实出行前开展多轮实际试用，持续以真实采用体验反哺规则与提示词迭代。
\end{itemize}
